% Generated draft: requires mathematical and source review.
\documentclass[11pt]{article}
\usepackage[margin=25mm]{geometry}
\usepackage{fontspec}
\setmainfont{Latin Modern Roman}
\usepackage{amsmath,amssymb,mathtools,unicode-math}
\setmathfont{Latin Modern Math}
\usepackage{enumitem,needspace,xurl,hyperref}
\setlength{\parindent}{0pt}
\setlength{\parskip}{5pt}
\setlength{\emergencystretch}{3em}
\newcommand{\N}{\mathbb{N}}
\newcommand{\Z}{\mathbb{Z}}
\newcommand{\Q}{\mathbb{Q}}
\newcommand{\R}{\mathbb{R}}
\newcommand{\C}{\mathbb{C}}
\newcommand{\F}{\mathbb{F}}
\newcommand{\A}{\mathbb{A}}
\newcommand{\PP}{\mathbb{P}}
\newcommand{\E}{\mathbb{E}}
\newcommand{\eps}{\varepsilon}
\newcommand{\dd}{\,\mathrm{d}}
\begin{document}
\section*{UnsolvedMath ID: 30001278}
\textbf{Source problem code:} catalogue entry
\subsection{Definitions and mathematical statement}
\emph{Unresolved scope.} The supplied record omits the hypotheses on the base field, the normalization of the valuation, the meaning of the different notation, the lattice used for reduction, and some parameter conventions. The following is a conditional reconstruction under explicit conventional choices; these choices require comparison with the cited report before this can be treated as its precise formulation.

Assume that \(p>2\) is prime, that \(K/\Q_p\) is a finite extension, and that \(G_K=\operatorname{Gal}(\overline K/K)\). Normalize \(v_K\) by \(v_K(K^\times)=\Z\), and put \(e=v_K(p)\). Assume that \(n,r\in\Z\) with \(n\geq1\) and \(r\geq1\). Let \(V\) be a finite-dimensional crystalline \(\Q_p\)-representation of \(G_K\) whose Hodge--Tate weights lie in \([0,r]\), with the sign convention for these weights still to be checked against the source. Choose a \(G_K\)-stable \(\Z_p\)-lattice \(T\subset V\), and set
\[
T_n=T/p^nT,\qquad
\rho_n:G_K\longrightarrow\operatorname{Aut}_{\Z/p^n\Z}(T_n),\qquad
L_n=\overline K^{\ker(\rho_n)}.
\]
Thus \(L_n/K\) is a finite Galois extension. Interpret \(\mathcal D_{L_n/K}\) as its different ideal. If \(v_{L_n}\) is the integer-normalized valuation on \(L_n\), use the convention
\[
v_K(\mathcal D_{L_n/K})
=\frac{v_{L_n}(\mathcal D_{L_n/K})}{e(L_n/K)},
\]
where \(e(L_n/K)\) is the relative ramification index.

For this conditional reconstruction, interpret the source's \(\mathbb N\) as \(\Z_{\geq0}\), and define \(\alpha'\) and \(\beta'\) by
\[
\frac{r}{p-1}=p^{\alpha'}\beta',\qquad
\alpha'\in\Z_{\geq0},\qquad p^{-1}<\beta'\leq1.
\]
The supplied proposed assertion is that, for every choice of the data satisfying these hypotheses,
\[
v_K(\mathcal D_{L_n/K})
\leq 1+e\bigl(n+\alpha'+\beta'\bigr)-p^{-(n+\alpha')}.
\]
The original wording asks to replace \(\alpha,\beta\) in the ramification bound of Theorem 2 by \(\alpha',\beta'\) defined above and prove that the resulting bound remains valid. Theorem 2 is not reproduced in the supplied record, so its complete scope and its agreement with this reconstruction remain unresolved. In particular, positivity and integrality of \(r\), the inclusion of zero in \(\mathbb N\), and the treatment of weight bound \(r=0\) are not supplied; the displayed decomposition cannot hold when \(r=0\).
\subsection{Short English statement}
The proposed bound controls the different of the finite extension detected by a crystalline representation modulo \(p^n\), using parameters determined by \(r/(p-1)\). The displayed formulation uses explicit provisional local-field and lattice conventions because the supplied record omits them. The record labels the problem open; this draft does not independently verify that status or establish the inequality.
\subsection{Sources}
\begin{itemize}
\item Supplied local UnsolvedMath registry record, permanent ID \(30001278\), title \emph{Sharper Ramification Bounds for Local Galois Representations}, source code \texttt{catalogue entry}, in the Oberwolfach Reports Open Problems collection. The record gives proposed year 2009 and supplies no individual proposer.
\item \emph{Algebraische Zahlentheorie} (2009), Oberwolfach Reports, pp. 1667--1730, DOI \texttt{10.4171/owr/2009/30}: \url{https://doi.org/10.4171/owr/2009/30}. The supplied background describes these reports as records of mathematical workshops and their open-problem discussions. The original problem refers to Theorem 2. The registry labels the statement \texttt{corrected\_verified} and reports that it was checked against Conjecture 3 in the official report, restoring the representation, field, different, and full sharpened ramification bound. These are supplied verification claims, not independent verification by this draft.
\item Supplied literature reference labelled ``Abrashkin on crystalline representations'': \url{https://doi.org/10.1070/IM1990v035n03ABEH000715}. No further bibliographic details or supporting theorem statement are supplied.
\item Literature assessment dated exactly \texttt{2026-08-21}: the supplied curation labels the problem open, says that the report describes the extension of the crystalline bounds as unclear, and reports finding no general theorem resolving this extension. Its listed sources are the OWR report and the Abrashkin reference above. This is supplied dated curation derived from the OWR material, not independently verified literature research; time-sensitive status and cited claims require independent checking.
\end{itemize}
\end{document}
